\section*{43. Idealdifferentiation und Differente.}

\begin{center}
\emph{Journal f. d. reine u. angew. Math. 188 (1950), S. 1--21}\\[0.75ex]
Von \emph{Emmy Noether} \(\dagger\)\srcfnmark{1)}.
\end{center}

\srcfntext{1)}{Die vorliegende von E. Noether nachgelassene, im Winter 1927/28 niedergeschriebene Arbeit wurde dem Herausgeber freundlicherweise von H. Grell zur Verfügung gestellt. Sie gibt eine ausführliche Darstellung des auf der D. M. V.-Tagung in Prag 1929 gehaltenen Vortrags gleichen Titels; siehe Jahresbericht D. M. V. 39 (1930), S. 17 kursiv. Ab § 6, 3 sollte das Manuskript offenbar noch genauer überarbeitet werden. Der Kenner wird jedoch die hier bis auf kleine Abrundungen unverändert gelassene skizzenhafte Darstellung ohne weiteres verstehen.}

\subsection*{Einleitung.}

Der Hauptsatz der Verzweigungstheorie des algebraischen Zahlkörpers sagt bekanntlich aus, daß die Differente (das Verzweigungsideal) mindestens durch die $(\varrho-1)$-te Potenz eines Primideals teilbar ist, das zur $\varrho$-ten Potenz in seiner Primzahl $p$ aufgeht, und zwar genau durch die $(\varrho-1)$-te Potenz, wenn $\varrho$ nicht durch $p$ teilbar ist, durch eine höhere, wenn $\varrho$ durch $p$ teilbar ist.

Dieser Satz ist das Analogon dazu, daß der Differentialquotient $f'(x)$ eines Polynoms $f(x)$ mindestens durch die $(\varrho-1)$-te Potenz eines Linearfaktors teilbar ist, der in $f(x)$ zur $\varrho$-ten Potenz aufgeht, und zwar genau durch die $(\varrho-1)$-te Potenz, wenn $\varrho$ nicht durch die Charakteristik des Koeffizientenbereichs teilbar ist, durch eine höhere, wenn $\varrho$ durch diese Charakteristik teilbar ist.

Ich zeige im folgenden, daß es sich hier um mehr als eine formale Analogie handelt. \emph{Die Differente läßt sich auffassen als Differentialquotient eines definierenden Ideals des Zahlkörpers $K$ in einem zugeordneten ganzzahligen Polynombereich von $x_1,\ldots,x_n$, der Differentialquotient genommen an der Stelle $x=\omega$, also für $x_1=\omega_1,\ldots,x_n=\omega_n$, wo $\omega_1,\ldots,\omega_n$ eine Modulbasis des Systems der ganzen Zahlen aus $K$ --- der Hauptordnung $\mathfrak o$ --- bedeutet.} Dabei entspricht das definierende Ideal $\mathfrak M$ der Gesamtheit der Relationen zwischen den $\omega$, besteht also aus allen ganzzahligen Polynomen $f(x)$, für die $f(\omega)$ verschwindet. Der Differentialquotient $\mathfrak M'[x\to\omega]$ aber ist definiert als Differenzenquotient, genommen an der Stelle $x=\omega$.

Man betrachte nämlich --- wegen $f(\omega)=0$ --- das Ideal $\mathfrak M$ als aus allen Differenzen $f(x)-f(\omega)$ bestehend, bilde weiter das Differenzenideal $\mathfrak B=(x_1-\omega_1,\ldots,x_n-\omega_n)$ und den Differenzenquotient $\mathfrak A=\mathfrak M:\mathfrak B$, wobei die Quotientenbildung den üblichen Idealquotienten bedeutet, genommen im Polynombereich der $x$ mit ganzen Zahlen aus $K$, also Zahlen aus $\mathfrak o$, als Koeffizienten. Dann wird die Differente $\mathfrak D$ definiert durch
\[
\mathfrak D=\mathfrak A[x\to\omega].
\]
Die Erweiterung des Koeffizientenbereichs ist notwendig, damit Differenzenideal und Differenzenquotient überhaupt definiert sind. Es handelt sich also um die direkte Verallgemeinerung des Differentialquotienten eines Polynoms einer Unbestimmten $x$
\[
f'(\xi)=(f(x)-f(\xi)):(x-\xi),
\]
genommen an der Stelle $x=\xi$, wobei auch hier der Polynombereich durch Adjunktion von $\xi$ erweitert werden muß, damit die Differenzen definiert sind. Tatsächlich geht auch der Idealdifferentialquotient in den Polynomdifferentialquotienten über, sobald die Basis der $\omega$ aus den Potenzen eines einzigen Elementes besteht.

Die hier gegebene Definition der Differente schließt somit direkt an Bekanntes an, sie führt aber nichtinvariante Zwischenglieder ein, insofern, als das Ideal $\mathfrak M$ von der gewählten Basis abhängt. Eine gleichbedeutende, \emph{vollständig invariante Definition} erhält man, wenn man statt des ganzzahligen Polynombereichs einen zur Hauptordnung $\mathfrak o$ isomorphen Ring $\mathfrak O$ einführt --- der dem Restklassenring nach $\mathfrak M$ isomorph wird --- und dessen Koeffizientenbereich durch $\mathfrak o$ zu $\mathfrak O_{\mathfrak o}$ erweitert.\srcfn{2)}{Die Erweiterung des Koeffizientenbereichs --- direkte Produktbildung --- wird in § 1 auch für nichtkommutative Ringe behandelt, allgemeiner als für die Zwecke dieser Arbeit erforderlich. Dafür genügt § 2 in der Spezialisierung auf kommutative Ringe und endliche Basis.} Das Differenzenideal $\mathfrak B$ wird hier aus allen Differenzen $(x-\xi)$ abgeleitet, wo $x$ und $\xi$ sich vermöge des Isomorphismus von $\mathfrak O$ zu $\mathfrak o$ entsprechen; der Differenzenquotient $\mathfrak A$ wird gleich dem Quotienten des Nullideals durch $\mathfrak B$; die Differente entsteht aus $\mathfrak A$, indem jeweils $x$ durch das isomorphe $\xi$ ersetzt wird.

Diese invariante Definition wird im folgenden an die Spitze gestellt, und zwar wird so direkt die Differente eines beliebigen kommutativen Ringes in bezug auf einen Unterring definiert, wobei nur gewisse Voraussetzungen in bezug auf die Möglichkeit der Koeffizientenerweiterung erfüllt sein müssen, was z. B. immer der Fall ist bei Existenz einer Modulbasis, eventuell unter Übergang zu gewissen Quotientenringen. Damit ist dann insbesondere auch die Differente einer beliebigen Ordnung des Zahlkörpers definiert, ebenso die Differente des Restklassenringes einer Ordnung etwa nach einer Primzahlpotenz; weiter auch die Relativdifferente und die Differente im Fall der algebraischen Funktionen von mehreren Unbestimmten. Die Übereinstimmung der Differentialdefinition der Differente mit der üblichen folgt aus einer genauen, auf der direkten Summenzerlegung beruhenden Strukturuntersuchung des dem Zahlkörper durch direkte Produktbildung zugeordneten galoisschen Erweiterungsringes\srcfn{3)}{Es handelt sich hier um eine Weiterführung der meiner Diskriminantenarbeit zugrunde liegenden Überlegungen, wobei aber die Kenntnis dieser Arbeit nicht vorausgesetzt wird. Vgl. E. Noether, Der Diskriminantensatz für die Ordnungen eines algebraischen Zahl- und Funktionenkörpers, J. Math. 157 (1927), 82--104.} --- eine Untersuchung, die im Spezialfall des Restklassenringes nach einem Polynom einer Unbestimmten auf eine Analyse der Lagrangeschen Interpolationsformel hinauskommt. Genauer verstehe ich darunter folgendes:

Es werde ein dem Zahlkörper $K$ isomorpher Ring $\mathfrak K$ eingeführt, der $\mathfrak O$ umfaßt, und sein Koeffizientenbereich --- die rationalen Zahlen --- durch den galoisschen Körper $F$ von $K$ erweitert. Dieses erweiterte System $\mathfrak K_F$ wird direkte Summe von $n$ einfachen Komponenten, die den $n$ Isomorphismen (Übergang zu den konjugierten Körpern) von $K$ entsprechen und Erweiterungen des Differenzenquotienten $\mathfrak A$ und seiner Konjugierten sind; zugleich wird das Nullideal von $\mathfrak K_F$ Durchschnitt der $n$ aus dem Differenzenideal und seinen Konjugierten abgeleiteten Ideale. Die einfachen Komponenten werden von der Form $Fe^{(i)}$, wo $e^{(i)}$ die Komponente der Einheit bedeutet; der Differenzenquotient wird --- da $e^{(i)}$ für $x=\xi$ in die Einheit übergeht --- gleich $\mathfrak d e^{(i)}$, unter $\mathfrak d$ die Differente von $\mathfrak o$ verstanden. Die Differente $\mathfrak d$ besteht aus allen Elementen des Körpers, die mit $e^{(i)}$ multipliziert in $\mathfrak O_{\mathfrak o}$ liegen, dem direkten Produkt der Ordnung mit sich selbst. Entsprechendes gilt für die Konjugierten.

Um von da aus die Definition durch den komplementären Modul zu gewinnen, ist zu beachten, daß $e^{(i)}$ eine Linearform in der den $\omega$ entsprechenden Basis $x_1,\ldots,x_n$ von $\mathfrak O$ bzw. $\mathfrak K$ wird; die Koeffizienten der $x$ in den $e^{(i)}$ aber ergeben gerade eine Basis des komplementären Moduls von $\mathfrak o$ (und seinen Konjugierten). Aus der Tatsache, daß der Differenzenquotient $\mathfrak A$ Koeffizienten aus $\mathfrak o$ hat, folgt sofort, daß die Differente gleich dem Quotienten von $\mathfrak o$ durch seinen Komplementärmodul wird, wodurch die Dedekindsche Definition gewonnen ist. Die angegebenen Betrachtungen sind nicht auf die Hauptordnung beschränkt; sie charakterisieren die Differente jeder Ordnung als Quotient der Ordnung durch ihren Komplementärmodul.

Die Übereinstimmung der Differentialdefinition mit der Definition durch die Fundamentalgleichung --- im Fall der Hauptordnung --- folgt aus der oben erwähnten Tatsache, daß der Idealdifferentialquotient in den Polynomdifferentialquotienten übergeht, sobald die Basis aus den Potenzen eines Elementes, hier der Fundamentalform, besteht. Damit ist dann auch der begriffliche Zusammenhang dieser letzteren Definition mit der Dedekindschen gegeben.

Unter Zugrundelegung der Fundamentalgleichung --- im Fall der Hauptordnung --- folgt bekanntlich direkt der am Anfang erwähnte Hauptsatz. Zu einer genauen Exponentenbestimmung führt die Betrachtung der Differente $\mathfrak d[p^t]$ des Restklassenringes der Hauptordnung nach $p^t$ mit genügend hohem $t$ --- es genügt $t=n$, was für Körper zweiten Grades auch notwendig ist --- wobei es wesentlich ist, daß die Differente $\mathfrak d$, betrachtet mod $p^t$, genau in $\mathfrak d[p^t]$ übergeht. Hier kann man sich wieder --- da der direkten Summe des Ringes die direkte Summe der Differenten entspricht --- auf den Restklassenring nach $\mathfrak p^{t\varrho}$ beschränken. Dieser letztere aber wird isomorph dem Restklassenring nach einem Polynom in einer Unbestimmten mit Koeffizienten mod $p^t$; und die Differentiation dieses Polynoms ergibt --- wie Ö. Ore, der auf anderem Wege zu einem solchen Polynom kam, gezeigt hat\srcfn{4)}{Anmerkung bei der Herausgabe: Dieses Zitat ist im Manuskript nicht ausgefüllt. Siehe die Hinweise in dem nur skizzierten § 7 der Arbeit.} --- den genauen Überblick über die überhaupt möglichen Exponenten mit Hilfe der Supplementzahlen.

Schließlich ergibt die Strukturuntersuchung noch den bekannten Satz, daß die Differente eines Oberkörpers von $K$ gleich dem Produkt der Relativdifferente mit der Differente von $K$ wird, und zwar als Folge der Tatsache, daß dieselbe Multiplikationseigenschaft für die Einheiten der direkten Summenzerlegung und damit für die Komplementärmoduln erfüllt ist.

Die angegebenen Tatsachen gelten überhaupt alle für Relativdifferenten durch Übergang zu gewissen Quotientenringen genau so, wie ich das in der Diskriminantenarbeit für Relativdiskriminanten ausgeführt habe. Es bleibt weiter --- durch Übergang zum Funktionalbereich --- alles erhalten für algebraische Funktionenkörper mit Koeffizientenkörper der Charakteristik Null, während bei Charakteristik $p$ und bei ganzzahligen algebraischen Funktionen zwar die Struktureigenschaften erhalten bleiben (im wesentlichen sogar ohne Übergang zum Funktionalbereich), die Verzweigungstheorie aber wegen des Auftretens von inseparablen Erweiterungskörpern sich ändert, was hier nicht mehr ausgeführt wird. Als wesentlichstes Ergebnis möchte ich die oben skizzierten vollständig invarianten, auch für algebraische Funktionen mehrerer Unbestimmten gültigen Struktursätze (§ 5 und § 6) bezeichnen; der Übergang zum Komplementärmodul bedeutet schon ein Auflösen in Koeffizientenidentitäten und ist nur durchgeführt, um Bekanntes einzuordnen.
% Source: Paper 43, printed pp. 4--6 (journal scan pp. 693--695)
\subsection*{§ 1. Erweiterung des Koeffizientenbereichs eines Ringes oder direkte Produktbildung. Definierendes Ideal\srcfnmark{5)}}
\srcfntext{5)}{Für kritische Bemerkungen zu diesem Paragraphen bin ich B. L. van der Waerden zu Dank verpflichtet.}

In diesem Paragraphen werden allgemeine Ringe zugrunde gelegt, für die das kommutative Gesetz der Multiplikation nicht vorausgesetzt wird. Dagegen wird der Einfachheit halber die Existenz des Einheitselementes vorausgesetzt, und zwar, sofern nicht ausdrücklich das Gegenteil vermerkt, in der ganzen Arbeit.

\paragraph{1. Definition des direkten Produktes.}
Ein Ring $\mathfrak O\times\mathfrak o$ oder $\mathfrak O_{\mathfrak o}$ heißt \emph{direktes Produkt} der Ringe $\mathfrak O$ und $\mathfrak o$ in bezug auf $\mathfrak h$ --- oder auch \emph{aus $\mathfrak O$ durch Erweiterung des Koeffizientenbereichs $\mathfrak h$ zu $\mathfrak o$ entstanden}\srcfn{6)}{Die zweite im folgenden meist gebrauchte Schreibweise soll an die Erweiterung des Koeffizientenbereichs erinnern; die erste ist die übliche Bezeichnung des direkten Produktes.} ---, wenn die folgenden Bedingungen erfüllt sind:

\begin{enumerate}
\item[(1)] $\mathfrak O_{\mathfrak o}$ enthält $\mathfrak O$ und $\mathfrak o$ als Unterring und wird gleich dem Produkt dieser Ringe (also gleich dem Durchschnitt aller $\mathfrak O$ und $\mathfrak o$ umfassenden Ringe in $\mathfrak O_{\mathfrak o}$).

\item[(2)] Der Durchschnitt von $\mathfrak O$ und $\mathfrak o$ ist durch $\mathfrak h$ gegeben, wo $\mathfrak h$ das Einheitselement von $\mathfrak O_{\mathfrak o}$ enthält.

\item[(3)] $\mathfrak O$ und $\mathfrak o$ sind elementweise vertauschbar; somit besteht $\mathfrak O_{\mathfrak o}$ aus der Gesamtheit der bilinearen Verbindungen
\[
x_{i_1}y_{i_1}+\cdots+x_{i_r}y_{i_r},
\]
wo die $x$ bzw. $y$ alle Systeme von je endlich vielen Elementen aus $\mathfrak O$ bzw. $\mathfrak o$ durchlaufen, und $\mathfrak h$ ist im Zentrum von $\mathfrak O$ und $\mathfrak o$ und damit von $\mathfrak O_{\mathfrak o}$ gelegen.

\item[(4)] Zwei Elemente aus $\mathfrak O_{\mathfrak o}$ sind nur dann (und natürlich stets dann) gleich, wenn sie auf mindestens eine Art formal gleich werden vermöge der in $\mathfrak O$ einerseits und in $\mathfrak o$ anderseits bestehenden Relationen. Darunter ist das folgende zu verstehen: Gilt
\[
x_1y_1+\cdots+x_ny_n=0
\quad\hbox{in }\mathfrak O_{\mathfrak o},
\]
so läßt sich dieser Ausdruck durch Zufügen von Summen formal verschwindender Verbindungen $(\gamma h)\delta-\gamma(h\delta)$, wo $h$ in $\mathfrak h$, $\gamma$ in $\mathfrak O$ und $\delta$ in $\mathfrak o$, umformen in
\[
(z_1\alpha_1+\cdots+z_r\alpha_r)+(t_1\beta_1+\cdots+t_s\beta_s)
\]
mit $z_1=0,\ldots,z_r=0$; $\beta_1=0,\ldots,\beta_s=0$, wo $\alpha,\beta$ in $\mathfrak o$ und $z,t$ in $\mathfrak O$, so daß also $z=0$ Relationen zwischen den $x,\gamma,h$ in $\mathfrak O$ und $\beta=0$ Relationen zwischen den $y,\delta,h$ in $\mathfrak o$ bedeuten.\srcfn{7)}{Setzt man die Existenz des Einheitselementes nicht voraus, so treten in (3) und (4) noch lineare Zusatzglieder auf.}
\end{enumerate}

Zwei Elemente aus $\mathfrak O_{\mathfrak o}$ sind nur dann gleich, wenn ihre Differenz auf die angegebene Art verschwindet.

Durch die Forderung des direkten Produktes sind also die Gleichheits- und Verknüpfungsbeziehungen (Addition und Multiplikation) eindeutig durch diejenigen der Faktoren festgelegt. Daraus folgt: Sind $\mathfrak O,\mathfrak o$ und ihr Durchschnitt $\mathfrak h$ isomorph zu $\overline{\mathfrak O},\overline{\mathfrak o}$ und $\overline{\mathfrak h}$, und existiert das direkte Produkt $\mathfrak O\times\mathfrak o$, so existiert auch $\overline{\mathfrak O}\times\overline{\mathfrak o}$ und wird zu $\mathfrak O\times\mathfrak o$ isomorph.

Jeder Ring $\mathfrak S$, der gleich dem Produkt von $\mathfrak O$ und $\mathfrak o$ wird, derart, daß $\mathfrak O$ und $\mathfrak o$ elementweise vertauschbar sind, wird homomorphes Bild von $\mathfrak O\times\mathfrak o$, wobei die Abbildung von $\mathfrak O$ zu $\overline{\mathfrak O}$ bzw. $\mathfrak o$ zu $\overline{\mathfrak o}$ isomorph bleibt. Denn die Gleichheit in $\mathfrak O\times\mathfrak o$ wird durch diejenige in $\mathfrak S$ umfaßt und stimmt für $\mathfrak O$ bzw. $\mathfrak o$ mit der von $\mathfrak O$ bzw. $\mathfrak o$ in $\mathfrak O\times\mathfrak o$ überein.\srcfn{8)}{Vergleiche auch § 1, 4.}

\paragraph{2. Notwendige und hinreichende Bedingungen für die Existenz des direkten Produkts in bezug auf einen Unterring.}
Sind $\mathfrak O$ und $\mathfrak o$ Ringe, deren Durchschnitt durch den im Zentrum beider gelegenen Ring $\mathfrak h$ gegeben ist, so braucht das direkte Produkt in bezug auf $\mathfrak h$ nicht zu existieren, wie das folgende Beispiel zeigt: Sei $\mathfrak h$ der Ring der ganzen rationalen Zahlen, $\mathfrak O=\mathfrak h[x]$ mit $1-2x=0$, und $\mathfrak o=\mathfrak h[y]$ mit $1-2y=0$, wo $y$ ein neues Symbol bedeutet. Es werden also $\mathfrak O$ und $\mathfrak o$ in bezug auf $\mathfrak h$ äquivalente Ringe, deren mengentheoretischer Durchschnitt durch $\mathfrak h$ gegeben ist, während zwischen den Symbolen $x$ und $y$ keine Verknüpfungen gegeben sind. Es gibt keinen $\mathfrak O$ und $\mathfrak o$ umfassenden Ring $R$ derart, daß in $R$ --- wo also jetzt Verknüpfungen zwischen $x$ und $y$ bestehen --- der Durchschnitt durch $\mathfrak h$ gegeben ist. Denn in $R$ folgt:
\[
0=(1-2x)y-x(1-2y)=y-x,
\]
wobei Vertauschbarkeit von $x$ und $y$ sogar nicht benutzt wird.\srcfn{9)}{Der in einem festen Ring gelegene Quotientenring eines kommutativen Ringes --- hier Adjunktion von $\frac12$ --- ist eben eindeutig. Es handelt sich um das Analogon dazu, daß zwei äquivalente und verschiedene galoissche Körpererweiterungen sich nicht in einen gemeinsamen Umfassungskörper einbetten lassen. Hier gibt es --- durch direkte Produktbildung --- allerdings noch Einbettung in einen Ring mit Nullteilern.}

Die \emph{notwendige und hinreichende Bedingung für die Existenz des direkten Produktes} ist die folgende:

\emph{Es muß mindestens einen $\mathfrak O$ und $\mathfrak o$ umfassenden Ring $R$ geben derart, daß der Durchschnitt von $\mathfrak O$ und $\mathfrak o$ durch $\mathfrak h$ gegeben ist und daß $\mathfrak O$ und $\mathfrak o$ elementweise vertauschbar werden.}

Existiert nämlich das direkte Produkt, so ist es selbst ein solcher Ring $R$. Um umgekehrt das direkte Produkt zu konstruieren, darf vorausgesetzt werden --- diese Voraussetzung wird im folgenden stets gemacht ---, daß zwischen den nicht zu $\mathfrak h$ gehörigen Elementen von $\mathfrak O$ und $\mathfrak o$ keine Verknüpfungen bestehen, da dies durch Übergang zu äquivalenter Erweiterung von $\mathfrak h$ --- Einführung anderer Symbole --- stets erreichbar ist.\srcfn{9a)}{Vgl. § 2, 1 und Fußnote 12).} Es enthalte weiter $\mathfrak h$ das Einheitselement $e$ von $\mathfrak O$ und $\mathfrak o$ (vgl. dazu aber Anmerkung 7)). Man definiere jetzt $\mathfrak T$ als Menge aller bilinearen Verbindungen
\[
x_{i_1}y_{i_1}+\cdots+x_{i_r}y_{i_r}
 =y_{i_1}x_{i_1}+\cdots+y_{i_r}x_{i_r},
\]
wo die $x$ und $y$ alle Systeme von je endlich vielen Elementen aus $\mathfrak O$ bzw. $\mathfrak o$ durchlaufen und die $x$ und $y$ also in $\mathfrak T$ vertauschbar werden. Man setze zwei solche Verbindungen dann und nur dann gleich, wenn sie vermöge der Relationen in $\mathfrak O$ und $\mathfrak o$ formal gleich werden in dem unter 1, (4) präzisierten Sinn. Man definiere weiter Summe und Produkt von
\[
\alpha_{i_1}x_{i_1}+\cdots+\alpha_{i_r}x_{i_r}
\quad\hbox{und}\quad
\beta_{i_1}x_{i_1}+\cdots+\beta_{i_r}x_{i_r},
\]
wo unter den $\alpha,\beta$ auch Nullen auftreten können, durch
\[
(\alpha_{i_1}+\beta_{i_1})x_{i_1}+\cdots+(\alpha_{i_r}+\beta_{i_r})x_{i_r}
\quad\hbox{und}\quad
\sum_{\lambda,\mu}\alpha_{i_\lambda}\beta_{i_\mu}x_{i_\lambda}x_{i_\mu}.
\]
Diese Definition ist eindeutig im Sinn der Gleichheit, da jede Nullrelation (im Sinne von 1, (4)) durch Rechts- oder Linksmultiplikation mit einem Element aus $\mathfrak o$ oder $\mathfrak O$ wieder in eine solche übergeht und dasselbe für Summe und Differenz dieser Nullrelationen gilt. Das System $\mathfrak T$ bildet einen $\mathfrak O$ und $\mathfrak o$ umfassenden Ring, der in bezug auf Gleichheit, Addition und Multiplikation den Bedingungen des direkten Produktes genügt. $\mathfrak T$ genügt aber auch der Durchschnittsbedingung als Folge der allgemeinen Einbettungsvoraussetzung. Sei nämlich $\mathfrak S$ das Produkt von $\mathfrak O$ und $\mathfrak o$ in $R$, dann wird $\mathfrak S$ nach 1 homomorphes Bild von $\mathfrak T$, isomorph in bezug auf $\mathfrak O$ und $\mathfrak o$. Eine Relation $x=y$ in $\mathfrak T$, wo $x$ in $\mathfrak O$ und $y$ in $\mathfrak o$, zieht somit dieselbe Relation in $\mathfrak S$ nach sich; also liegen in $\mathfrak S$ nach Voraussetzung beide Elemente in $\mathfrak h$, und wegen des Isomorphismus in bezug auf $\mathfrak O$ und $\mathfrak o$ einzeln gilt dasselbe für $\mathfrak T$. Die Einbettungsvoraussetzung ist also notwendig und hinreichend.

\paragraph{3. Definierendes Ideal eines Ringes in bezug auf einen Unterring.}
Ein System $S$ von Elementen $\ldots,z,\ldots$ aus $\mathfrak O$ heißt \emph{Erzeugendensystem} von $\mathfrak O$ in bezug auf $\mathfrak h$, wenn $\mathfrak O$ gleich dem durch $\mathfrak h$ und $S$ in $\mathfrak O$ erzeugten Ring wird --- also gleich dem Durchschnitt aller $\mathfrak h$ und $S$ umfassenden Ringe aus $\mathfrak O$; in Zeichen $\mathfrak O=\mathfrak h[S]$. Es sei wieder $\mathfrak h$ als Unterring des Zentrums von $\mathfrak O$ vorausgesetzt; es wird dann $\mathfrak O$ homomorphes Bild des nichtkommutativen Polynombereichs $\mathfrak h[\ldots Z\ldots]$, wo $Z$ Unbestimmte bedeuten, und das System der Unbestimmten von gleicher Mächtigkeit ist wie $S$. Unter dem nichtkommutativen Polynombereich wird dabei das System aller endlichen Summen $\sum h_iP_i=\sum P_i h_i$ verstanden, wo die $P_i$ Potenzprodukte einschließlich der dem Exponenten $0$ entsprechenden Einheit sind:
\[
P_i=Z_{i_1}^{\varrho_1}Z_{i_2}^{\varrho_2}\cdots Z_{i_n}^{\varrho_n}
\]
(wobei natürlich die gleichen Indizes sich wiederholen können), und wo die Gleichheit als Koeffizientengleichheit in bezug auf die verschiedenen Potenzprodukte definiert ist. Aus dem Homomorphismus folgt, daß $\mathfrak O$ isomorph wird dem Restklassenring $\mathfrak h[\ldots Z\ldots]/\mathfrak M$, wo das zweiseitige Ideal $\mathfrak M$ aus allen Polynomen $F(Z)$ (Elementen aus $\mathfrak h[\ldots Z\ldots]$) besteht derart, daß $F(z)$ in $\mathfrak O$ verschwindet --- wo $Z$ und $z$ sich vermöge des Homomorphismus entsprechen. \emph{Das Ideal $\mathfrak M$ heißt ein definierendes Ideal in bezug auf $\mathfrak h$, und zwar das dem Erzeugendensystem der $z$ entsprechende Ideal.} Die Definition zeigt, daß isomorphe Ringe in bezug auf isomorph entsprechende Unterringe $\mathfrak h,\overline{\mathfrak h}$ und Erzeugendensysteme $S$ und $\overline S$ auch isomorphe definierende Ideale besitzen --- die insbesondere, wenn $\overline{\mathfrak h}$ mit $\mathfrak h$ zusammenfällt, auch als identisch angenommen werden dürfen.

Da $\mathfrak h$ im Zentrum von $\mathfrak O$ angenommen war, wird $\mathfrak O$ sicher kommutativ, wenn es ein Erzeugendensystem $S$ besitzt, das aus einem einzigen Element $z$ besteht. Wird $\mathfrak M$ noch Hauptideal, so heißt die einem Basispolynom $G(Z)$ von $\mathfrak M$ entsprechende Gleichung $G(z)=0$ eine definierende Gleichung von $\mathfrak O$ in bezug auf $\mathfrak h$.

\paragraph{4. Das definierende Ideal eines direkten Produktes.}
Es sei (wie unter 3) $\mathfrak O$ isomorph zu $\mathfrak h[\ldots Z\ldots]/\mathfrak M$, wo also $\mathfrak M$ definierendes Ideal. Neben $\mathfrak h[\ldots Z\ldots]$ werde der Polynombereich $\mathfrak o[\ldots Z\ldots]$ betrachtet, also das System aller endlichen Summen
\[
\sum \gamma_iP_i=\sum P_i\gamma_i
\qquad\hbox{mit }\gamma_i\hbox{ aus }\mathfrak o,
\]
die Gleichheit als Koeffizientengleichheit definiert. $\mathfrak o[Z]$ umfaßt $\mathfrak h[Z]$ und ist direktes Produkt von $\mathfrak h[\ldots Z\ldots]$ und $\mathfrak o$ in bezug auf $\mathfrak h$, anders ausgedrückt: aus $\mathfrak h[\ldots Z\ldots]$ durch Erweiterung des Koeffizientenbereichs hervorgegangen.\srcfn{10)}{Denn Gleichheits- und Verknüpfungsbedingungen sind erfüllt, und aus $\gamma=h+h_1P_1+\cdots+h_rP_r$ folgt nach der Gleichheit in $\mathfrak o[\ldots Z\ldots]$, daß $\gamma=h$. Vgl. dazu § 2, 2. Die Tatsache, daß $\mathfrak o[\ldots Z\ldots]$ direktes Produkt ist, wird in dieser Nummer nicht benutzt.} Es bedeute $\mathfrak M_{\mathfrak o}$ das Erweiterungsideal von $\mathfrak M$ in bezug auf $\mathfrak o[\ldots Z\ldots]$ --- Durchschnitt aller $\mathfrak M$ umfassenden zweiseitigen Ideale aus $\mathfrak o[\ldots Z\ldots]$ ---, also das System aller linearen Verbindungen der Elemente aus $\mathfrak M$ mit Koeffizienten aus $\mathfrak o$, das wegen der Vertauschbarkeit von $\mathfrak o$ mit den $Z$ schon ein zweiseitiges Ideal bildet. \emph{Falls das direkte Produkt $\mathfrak O_{\mathfrak o}$ in bezug auf $\mathfrak h$ existiert, so wird $\mathfrak O_{\mathfrak o}$ isomorph zu $\mathfrak o[\ldots Z\ldots]/\mathfrak M_{\mathfrak o}$.} Denn $\mathfrak O_{\mathfrak o}$ wird homomorphes Bild von $\mathfrak o[\ldots Z\ldots]$, also isomorph dem Restklassenring nach dem definierenden Ideal von $\mathfrak O_{\mathfrak o}$ in bezug auf $\mathfrak o$, entsprechend dem Erzeugendensystem der $z$. Nach der Gleichheitsdefinition 1, (4) aber lassen sich die zwischen den $z$ in $\mathfrak O_{\mathfrak o}$ bestehenden Relationen als lineare Kombinationen mit Koeffizienten aus $\mathfrak o$ der in $\mathfrak O$ bestehenden schreiben; also wird $\mathfrak M_{\mathfrak o}$ definierendes Ideal.

\paragraph{Bemerkung.}
Dieselbe Überlegung ergibt noch eine symmetrische, aber umständlichere Darstellung des direkten Produktes und des zugehörigen definierenden Ideals. Es sei $\mathfrak o$ isomorph $\mathfrak h[\ldots Y\ldots]/\mathfrak N$, wo $Y$ von den $Z$ verschiedene Symbole bedeuten, so daß der Polynombereich $\mathfrak h[\ldots Y\ldots Z\ldots]$ definiert ist. Dann wird $\mathfrak O_{\mathfrak o}$ isomorph
\[
\mathfrak h[\ldots Y\ldots Z\ldots]/(\mathfrak M,\mathfrak N,\ldots,YZ-ZY,\ldots),
\]
wo das rechtsstehende definierende Ideal durch $\mathfrak M,\mathfrak N$ und alle Kombinationen $YZ-ZY$ erzeugt ist. (Letztere drücken die Vertauschbarkeit der Elemente von $\mathfrak o$ und $\mathfrak O$ aus, erstere die Tatsache, daß die Relationen in $\mathfrak O_{\mathfrak o}$ Folgen der in $\mathfrak O$ und $\mathfrak o$ einzeln sind.)
% Source: Paper 43, printed pp. 7--10 (journal scan pp. 696--699)
\subsection*{§ 2. Ringe mit unabhängiger Modulbasis. Konstruktion des direkten Produktes.}

\paragraph{1. Konstruktion des direkten Produktes.}
Ein System $T$ von Elementen $\ldots,t,\ldots$ aus $\mathfrak O$ heißt eine $\mathfrak h$-Modulbasis von $\mathfrak O$, wenn $\mathfrak O$ gleich dem aus $T$ in $\mathfrak O$ abgeleiteten $\mathfrak h$-Modul wird. Es besteht dann $\mathfrak O$ wegen der Existenz der Einheit in $\mathfrak h$ aus allen linearen Verbindungen $h_{i_1}t_{i_1}+\cdots+h_{i_r}t_{i_r}$. Die Modulbasis heißt unabhängig, wenn aus
\[
h_{i_1}t_{i_1}+\cdots+h_{i_r}t_{i_r}=0
\quad\hbox{stets}\quad
h_{i_1}=0,\ldots,h_{i_r}=0
\]
folgt. \emph{Besitzt $\mathfrak O$ eine unabhängige, das Einheitselement enthaltende Modulbasis, so existiert das direkte Produkt $\mathfrak O_{\mathfrak o}$ bei beliebigem Erweiterungsring $\mathfrak o$ von $\mathfrak h$, für den nur $\mathfrak h$ Unterring des Zentrums und dessen mengentheoretischer Durchschnitt $[\mathfrak O,\mathfrak o]$ durch $\mathfrak h$ gegeben ist.}\srcfn{11)}{Natürlich könnte symmetrisch dazu auch die Basisvoraussetzung für $\mathfrak o$ gemacht werden. Nach § 2, 2 wird --- ohne Beschränkung der Allgemeinheit; vgl. E. Noether, Der Diskriminantensatz für die Ordnungen eines algebraischen Zahl- oder Funktionenkörpers, dies. Journ. 157 (1927), 82--104, §§ 1, 2 --- vorausgesetzt, daß zwischen den nicht zu $\mathfrak h$ gehörigen Elementen aus $\mathfrak O$ und $\mathfrak o$ keine Verknüpfungen bestehen.} Es genügt nach § 1, 2, einen entsprechenden Einbettungsring $\mathfrak R$ zu konstruieren, der sich dann hier schon als das direkte Produkt erweisen wird. Es sei $\mathfrak R$ definiert als Gesamtheit aller linearen Verbindungen
\[
\gamma_{i_1}t_{i_1}+\cdots+\gamma_{i_r}t_{i_r}
=t_{i_1}\gamma_{i_1}+\cdots+t_{i_r}\gamma_{i_r},
\]
wobei die $t_{i_\nu}$ alle Systeme je endlich vieler Elemente der unabhängigen Basis $T$ durchlaufen, die $\gamma$ entsprechend aus $\mathfrak o$. Die Gleichheit sei als Koeffizientengleichheit in den $t_{i_\nu}$ definiert, Summe und Produkt nach den üblichen Rechenregeln, wie unter § 1, 2. Wegen der Basiseigenschaft ist das Produkt irgend zweier $t_{i_\nu}$ eine lineare Verbindung endlich vieler mit Koeffizienten aus $\mathfrak h$, es wird also $\mathfrak R$ ein Ring. $\mathfrak R$ umfaßt $\mathfrak O$, da $\mathfrak h$ Unterring von $\mathfrak o$ ist; und es umfaßt $\mathfrak o$, da $T$ das Einheitselement enthält. Da weiter die $\gamma$ sowohl mit den $h_{i_\nu}$ wie mit den $t_{i_\nu}$ elementweise vertauschbar sind, wird $\mathfrak O$ in $\mathfrak R$ mit $\mathfrak o$ elementweise vertauschbar. Schließlich ist in $\mathfrak R$ der Durchschnitt von $\mathfrak O$ und $\mathfrak o$ durch $\mathfrak h$ gegeben, denn eine Relation
\[
\gamma=\gamma e=h_1t_1+\cdots+h_rt_r
\]
ergibt nach der Gleichheitsdefinition --- da $e$ nach Voraussetzung Element der Basis ist ---, daß ein $t$ gleich $e$ ist, die übrigen verschwinden, und somit $\gamma=h$. Somit existiert das direkte Produkt; es wird gleich $\mathfrak R$, das gleich dem Produkt von $\mathfrak O$ und $\mathfrak o$ ist und der Vertauschbarkeits- und Gleichheitsbedingung genügt (da die Elemente der Basis unabhängig sind, geht diese Bedingung in Koeffizientengleichheit über). Diese Basis $T$ von $\mathfrak O$ wird zugleich $\mathfrak o$-Basis von $\mathfrak O_{\mathfrak o}$.

\paragraph{2. Erweiterungs- und Verengungsmoduln.}
Ist $\mathfrak B$ ein $\mathfrak h$-Modul aus $\mathfrak O$, so wird unter seiner Erweiterung $\mathfrak B_{\mathfrak o}$ der $\mathfrak o$-Modul aus $\mathfrak O_{\mathfrak o}$ verstanden, der Durchschnitt aller $\mathfrak B$ umfassenden $\mathfrak o$-Moduln aus $\mathfrak O_{\mathfrak o}$ ist. Ist $\mathfrak C$ ein $\mathfrak o$-Modul aus $\mathfrak O_{\mathfrak o}$, so ist sein Verengungsmodul $\mathfrak C_{\mathfrak h}$ in $\mathfrak O$ als Durchschnitt $[\mathfrak C,\mathfrak O]$ definiert. Ist $\mathfrak B$ zugleich Ideal in $\mathfrak O$, so wird $\mathfrak B_{\mathfrak o}$ Ideal in $\mathfrak O_{\mathfrak o}$; ist $\mathfrak C$ Ideal in $\mathfrak O_{\mathfrak o}$, so wird $[\mathfrak C,\mathfrak O]$ Ideal in $\mathfrak O$.

\paragraph{Satz.}
\emph{Besitzt ein Modul $\mathfrak B$ aus $\mathfrak O$ eine unabhängige $\mathfrak h$-Modulbasis, die sich zu einer unabhängigen $\mathfrak h$-Modulbasis von $\mathfrak O$ ergänzen läßt, so ist $\mathfrak B$ Verengung seines Erweiterungsmoduls $\mathfrak B_{\mathfrak o}$ in $\mathfrak O_{\mathfrak o}$.}

\paragraph{Beweis.}
Es bedeute $T$ eine unabhängige $\mathfrak h$-Modulbasis von $\mathfrak B$, die sich nach Voraussetzung durch ein System $\overline T$ von Elementen $\overline t$ aus $\mathfrak O$ zu einer unabhängigen $\mathfrak h$-Modulbasis von $\mathfrak O$ ergänzen läßt. Es wird $T$ dann zugleich unabhängige $\mathfrak o$-Basis von $\mathfrak B_{\mathfrak o}$, und $T$ und $\overline T$ zusammen eine solche von $\mathfrak O_{\mathfrak o}$.

Es ist $\mathfrak B$ enthalten im Durchschnitt $[\mathfrak B_{\mathfrak o},\mathfrak O]$. Sei umgekehrt $c$ ein Element dieses Verengungsmoduls. Als Element von $\mathfrak B_{\mathfrak o}$ läßt $c$ eine Darstellung durch Elemente von $T$ zu, die zugleich eine Darstellung durch die $\mathfrak o$-Basis $T,\overline T$ von $\mathfrak O_{\mathfrak o}$ bedeutet:
\[
c=\gamma_1t_1+\cdots+\gamma_rt_r
\qquad\hbox{mit }\gamma\hbox{ aus }\mathfrak o.
\]
Als Element von $\mathfrak O$ besitzt $c$ eine Darstellung durch die $\mathfrak h$-Basis $T,\overline T$ von $\mathfrak O$, also:
\[
c=h_{i_1}t_{i_1}+\cdots+h_{i_r}t_{i_r}+h_{j_1}\overline t_{j_1}+\cdots+h_{j_s}\overline t_{j_s},
\]
die ebenfalls zugleich, da $\mathfrak O$ Unterring von $\mathfrak O_{\mathfrak o}$ und $\mathfrak h$ Unterring von $\mathfrak o$, eine Darstellung durch die $\mathfrak o$-Basis $T,\overline T$ von $\mathfrak O_{\mathfrak o}$ bedeutet. Also müssen beide Darstellungen koeffizientengleich sein; die Koeffizienten der $\overline t$ verschwinden und
\[
c=h_1t_1+\cdots+h_rt_r
\]
zeigt, daß $c$ Element aus $\mathfrak B$ ist. Damit ist $\mathfrak B$ als Verengungsmodul erkannt.

\paragraph{Zusatzbemerkung.}
Der Beweis benutzt nur etwas geringere Voraussetzungen als die im Satz formulierten und beweist die übrigen; es gilt die verschärfte Fassung:

\emph{Es wird $\mathfrak B$ Verengungsmodul seiner Erweiterung $\mathfrak B_{\mathfrak o}$, wenn es in $\mathfrak O$ eine unabhängige $\mathfrak h$-Modulbasis $T,\overline T$ gibt, wo $T$ in $\mathfrak B$, derart, daß $T$ zugleich (unabhängige) $\mathfrak o$-Modulbasis von $\mathfrak B_{\mathfrak o}$ ist. Es erweist sich dann $T$ auch als unabhängige $\mathfrak h$-Modulbasis von $\mathfrak B$.}

\paragraph{3. Hinreichendes Kriterium für Existenz des direkten Produktes.}
Es kann vorkommen (z. B. bei der Relativdifferente), daß erst bei Erweiterung des Koeffizientenbereichs $\mathfrak h$ eine unabhängige Modulbasis existiert. Die Existenz des direkten Produktes kann sich dann entnehmen lassen aus dem hinreichenden Kriterium:

\emph{Gibt es einen Erweiterungsring $\mathfrak f$ von $\mathfrak h$ derart, daß die direkten Produkte $\mathfrak O_{\mathfrak f}$ und $\mathfrak o_{\mathfrak f}$ in bezug auf $\mathfrak h$ und weiter $\mathfrak O_{\mathfrak f}\times\mathfrak o_{\mathfrak f}$ in bezug auf $\mathfrak f$ existieren, so existiert $\mathfrak O\times\mathfrak o$ in bezug auf $\mathfrak h$.}

Es ist zu zeigen, daß $\mathfrak R=\mathfrak O_{\mathfrak f}\times\mathfrak o_{\mathfrak f}$ ein Einbettungsring im Sinn von § 1, 2 wird. In der Tat werden $\mathfrak O$ und $\mathfrak o$ als Unterringe von $\mathfrak O_{\mathfrak f}$ und $\mathfrak o_{\mathfrak f}$ elementweise vertauschbar; weiter gilt
\[
[\mathfrak O,\mathfrak o]\subseteq[\mathfrak O_{\mathfrak f},\mathfrak o_{\mathfrak f}]=\mathfrak f;
\]
ebenso $[\mathfrak O,\mathfrak o]\subseteq\mathfrak O$; also $[\mathfrak O,\mathfrak o]\subseteq[\mathfrak O,\mathfrak f]=\mathfrak h$, und somit $[\mathfrak O,\mathfrak o]=\mathfrak h$, womit $\mathfrak R$ als Einbettungsring nachgewiesen ist.

\paragraph{Bemerkung.}
Der Beweis benutzt die Durchschnittsvoraussetzung nur bei $\mathfrak O_{\mathfrak f}$, nicht bei $\mathfrak o_{\mathfrak f}$, welch letzterer Ring also auch einfach das Produkt von $\mathfrak o$ und $\mathfrak f$ sein kann.

\subsection*{§ 3. Differente eines Ringes nach einem Unterring. Differentialquotient eines definierenden Ideals.}

Von jetzt an seien alle auftretenden Ringe als kommutativ angenommen.\srcfn{12)}{Ohne diese Voraussetzung müßten Rechts- und Linksdifferente definiert und es müßte untersucht werden, wie weit diese unter sich übereinstimmen und wie weit sie im Fall von Ordnungen halbeinfacher Systeme mit bekannten Bildungen übereinstimmen (als Verallgemeinerung der Untersuchungen von § 4).} Weiter sei stets vorausgesetzt, daß die Ringe $\mathfrak O$ und $\mathfrak o$ isomorphe, und zwar äquivalente Erweiterungen des Unterringes $\mathfrak h$ seien (daß also die Abbildung von $\mathfrak O$ auf $\mathfrak o$ Fortsetzung der identischen von $\mathfrak h$ sei) und daß das direkte Produkt $\mathfrak O_{\mathfrak o}$ in bezug auf $\mathfrak h$ existiert.

\paragraph{1. Definition der Differente.}
Es sei vorausgesetzt, daß das direkte Produkt $\mathfrak O_{\mathfrak o}$ in bezug auf $\mathfrak h$ existiert. Es durchlaufe $x$ alle Elemente aus $\mathfrak O$, und $\xi$ bedeute jeweils das vermöge des Isomorphismus in $\mathfrak o$ entsprechende; es bezeichne $\mathfrak B$ das durch alle Differenzen $x-\xi$ erzeugte zweiseitige \emph{Differenzenideal} von $\mathfrak O_{\mathfrak o}$,
\[
\mathfrak B=\{\ldots,(x-\xi),\ldots\}.\srcfnmark{13)}
\]
\srcfntext{13)}{Diese Bezeichnung soll immer für aus Gesamtheiten abgeleitete Ideale angewandt werden.}
Der \emph{Differenzenquotient} $\mathfrak A$ ist definiert als Idealquotient des Nullideals durch $\mathfrak B$, besteht also aus allen Elementen $a$ aus $\mathfrak O_{\mathfrak o}$, für die $a\mathfrak B$ verschwindet,
\[
\mathfrak A=(0):\mathfrak B.
\]
Die \emph{Differente} $\mathfrak d$ von $\mathfrak o$ in bezug auf $\mathfrak h$ ist definiert als $\mathfrak A[x\to\xi]$, entsteht also aus $\mathfrak A$, indem in jedem $a=\alpha_1x_1+\cdots+\alpha_rx_r$ aus $\mathfrak A$ die $x$ durch die isomorph entsprechenden $\xi$ ersetzt werden. Entsprechend ist die Differente $\mathfrak D$ von $\mathfrak O$ in bezug auf $\mathfrak h$ als $\mathfrak A[\xi\to x]$ definiert; da $\mathfrak A$ symmetrisch in $x$ und $\xi$ definiert, entsprechen sich $\mathfrak d$ und $\mathfrak D$, als durch Vertauschen von $\xi$ und $x$ hervorgegangen, vermöge des Isomorphismus von $\mathfrak o$ zu $\mathfrak O$. Da $\mathfrak A$ Ideal in $\mathfrak O_{\mathfrak o}$, also insbesondere $\mathfrak o$-Modul und $\mathfrak O$-Modul, werden $\mathfrak d$ und $\mathfrak D$ Ideale in $\mathfrak o$ und $\mathfrak O$.

\paragraph{2. Differentialquotient eines definierenden Ideals.}
Es sei (§ 1, 3) $\mathfrak M$ definierendes Ideal von $\mathfrak O$ in bezug auf $\mathfrak h$, entsprechend dem Erzeugendensystem der $z$. Es wird also (§ 1, 3) $\mathfrak M$ auch definierendes Ideal von $\mathfrak o$ in bezug auf $\mathfrak h$, entsprechend dem isomorphen Erzeugendensystem der $\xi$; und es wird (§ 1, 4) $\mathfrak O_{\mathfrak o}$ isomorph $\mathfrak o[\ldots Z\ldots]/\mathfrak M_{\mathfrak o}$. Differenzenideal und Differenzenquotient sind jetzt im Polynombereich $\mathfrak o[\ldots Z\ldots]$ definiert durch
\[
\mathfrak Q=\{\ldots,(Z-\xi),\ldots\}
\]
und
\[
\mathfrak C=\mathfrak M_{\mathfrak o}:\mathfrak Q
=\{\ldots(F(Z)-F(\xi)),\ldots\}:\{\ldots,(Z-\xi),\ldots\}.
\]
Der \emph{Differentialquotient} $\mathfrak M'$ ist definiert als $\mathfrak C[\xi\to Z]$. Dabei ist zu beachten, daß diese Definition eindeutig ist, unabhängig von der \emph{Darstellung} der Elemente aus $\mathfrak o$ durch die Erzeugenden $\xi$. Denn da $\mathfrak C$ Idealquotient, umfaßt es $\mathfrak M_{\mathfrak o}$; zwei Darstellungen eines Elementes von $\mathfrak o$ unterscheiden sich durch ein Polynom $F(\xi)$, also nach der Einsetzung von $Z$ durch ein in $\mathfrak M$ und damit schon in $\mathfrak C$ liegendes $F(Z)$. Das Ideal $\mathfrak M'$ wird ein $\mathfrak M$ umfassendes Ideal in $\mathfrak h[\ldots Z\ldots]$. \emph{Der Differentialquotient $\mathfrak M'$ geht für $Z\to\xi$ in die Differente $\mathfrak d$ von $\mathfrak o$ über.} Denn vermöge einer homomorphen Abbildung entsprechen sich neben Idealen auch die Quotienten. Es wird also $\mathfrak A$ homomorphes Bild von $\mathfrak C$ vermöge des Homomorphismus zwischen $\mathfrak o[\ldots Z\ldots]$ und $\mathfrak O_{\mathfrak o}$, also des Entsprechens von $Z$ und $z$, da das gleiche für $\mathfrak B$ und das Nullideal in bezug auf $\mathfrak Q$ und $\mathfrak M_{\mathfrak o}$ gilt.\srcfn{13a)}{Anmerkung bei der Herausgabe: Es ist wichtig, daß $\mathfrak M_{\mathfrak o}$ die Urbildmenge von $(0)$ darstellt, um das Entsprechen des Quotienten schließen zu können.} Daraus folgt, daß der Substitution $\xi\to Z$ in $\mathfrak C$ vermöge des Homomorphismus die Substitution $\xi\to z$ in $\mathfrak A$ entspricht; es wird also die Differente $\mathfrak D$ von $\mathfrak O$ homomorphes Bild von $\mathfrak M'$ vermöge des Entsprechens von $Z$ und $z$. Also
\[
\mathfrak M'[Z\to z]=\mathfrak D
\qquad\hbox{und}\qquad
\mathfrak M'[Z\to\xi]=\mathfrak D[z\to\xi]=\mathfrak d.
\]

\paragraph{3. Differente im Fall einer definierenden Gleichung.}
Es besitze $\mathfrak O$ eine definierende Gleichung $G(z)=0$ (§ 1, 3, Schluß), wobei weiter noch \emph{vorausgesetzt} sei, daß der Koeffizient des höchsten Gliedes gleich der Einheit wird, also
\[
G(Z)=Z^n+h_1Z^{n-1}+\cdots+h_n
\]
mit Koeffizienten $h_i$ aus $\mathfrak h$. Dann ist der Zusammenhang mit dem Polynomdifferentialquotienten hergestellt durch:

\paragraph{Satz.}
\emph{Die Differente von $\mathfrak o$ bzw. $\mathfrak O$ in bezug auf $\mathfrak h$ ist definiert (d. h. das direkte Produkt $\mathfrak O_{\mathfrak o}$ in bezug auf $\mathfrak h$ existiert) und wird Hauptideal mit dem Basiselement $G'(\xi)$ bzw. $G'(z)$, wo $G'(Z)$ den Polynomdifferentialquotienten bedeutet. Der Idealdifferentialquotient $\mathfrak M'$ wird gleich $(G'(Z),G(Z))$.}

\paragraph{Beweis.}
Aus der Voraussetzung folgt zuerst, daß die Elemente $e,z,\ldots,z^{n-1}$ eine unabhängige Modulbasis von $\mathfrak O$ bilden. Denn da der höchste Koeffizient in $G(Z)$ die Einheit ist, bilden diese Potenzen tatsächlich eine $\mathfrak h$-Modulbasis; diese ist aber unabhängig, weil $\mathfrak M$ kein Polynom von niedrigerem als $n$-tem Grad enthalten kann. Der Grad eines Polynoms $K(Z)G(Z)$ ist nämlich gleich der Summe der Gradzahlen der Faktoren, wieder wegen der Voraussetzung über den höchsten Koeffizienten. Wegen der Existenz der Modulbasis existiert also das direkte Produkt $\mathfrak O_{\mathfrak o}$ in bezug auf $\mathfrak h$, und damit ist die Differente definiert. Das Weitere folgt aus der Betrachtung des Differenzenquotienten $\mathfrak C$ in $\mathfrak o[\ldots Z\ldots]$ (vgl. 2). Dieser Differenzenquotient wird Hauptideal mit einem Basispolynom
\[
\begin{aligned}
C(Z,\xi)={}&(Z^{n-1}+Z^{n-2}\xi+\cdots+\xi^{n-1})\\
&+h_1(Z^{n-2}+\cdots+\xi^{n-2})+\cdots+h_{n-1}.
\end{aligned}
\]
Denn $C$ entsteht durch formale Division von $G(Z)-G(\xi)$ durch $Z-\xi$ und gehört somit zu $\mathfrak C$, da $Z-\xi$ Basis von $\mathfrak Q$ wird. Ein beliebiges Polynom läßt sich mod $C$ auf ein solches $D$ reduzieren, das in $Z$ den Grad $n-1$ nicht erreicht; also erreicht das Produkt $D(Z-\xi)$ den Grad $n$ nicht und kann daher nur zu $\mathfrak M$ gehören, wenn es verschwindet. Aus $D(Z-\xi)=0$ folgt aber $D=0$, da der Koeffizient von $Z$ in $Z-\xi$ die Einheit ist. Also wird $C$ die Basis von $\mathfrak C$. Nun aber ergibt sich
\[
\mathfrak M'=(G'(Z),G(Z)),
\]
wie die Substitution $\xi\to Z$ in $C$ und seine Vielfachen (also $G(Z)=C(Z-\xi)$ und dessen Vielfache) zeigt, und damit kommt
\[
\mathfrak d=(G'(\xi))
\qquad\hbox{und}\qquad
\mathfrak D=(G'(z)).
\]
% Source: Paper 43, printed pp. 10--13 (journal scan pp. 699--702)
